\chapter{Clipping and overflow}
\label{ch:clipping}

\begin{aside}
A widget that doesn't fit shouldn't break the screen. Layout
gives every element a rectangle; clipping enforces it. The
question is what to do with what didn't fit.
\end{aside}

\section{The clip rectangle}

Each element receives a clip rectangle from its parent. The
painter respects it: any \code{put(x, y, cell)} outside the
clip is dropped, not drawn.

\begin{lstlisting}
void Painter::put(int x, int y, Cell c) {
    if (!clip_.contains(x, y)) return;
    back_(x, y) = c;
}
\end{lstlisting}

The check is one comparison and one branch. Modern CPUs do
this in roughly half a cycle.

\section{Overflow strategies}

A child too wide or tall for its parent can be handled in
five ways:

\begin{itemize}[leftmargin=*]
  \item \textbf{Truncate.} Drop what doesn't fit. Default.
  \item \textbf{Ellipsis.} Drop and append \code{...} to the
    last visible cell.
  \item \textbf{Scroll.} Keep all content, attach a viewport
    with scroll position.
  \item \textbf{Wrap.} Break onto a new row.
  \item \textbf{Resize parent.} Grow the parent. Usually
    impossible (the parent has its own constraints).
\end{itemize}

\maya{} exposes these as element modifiers:

\begin{lstlisting}
auto t = text("very long line of stuff")
       | overflow(Overflow::Ellipsis)
       | width(20);
\end{lstlisting}

\section{Ellipsis details}

The ellipsis is a single cell containing \code{U+2026}
HORIZONTAL ELLIPSIS (or three dots in ASCII fallback). The
last visible cell becomes the ellipsis; if the text is itself
shorter than the width, no ellipsis is drawn.

A subtle case: the cell before the ellipsis might be a wide
character whose trailing slot would now be visible without
its head. \maya{} extends the truncation by one to avoid
half-wide-cell artifacts.

\section{Wrap}

Wrap mode pushes overflow onto subsequent rows. The element
grows in height (subject to its height constraint).

The word-boundary detection is grapheme-aware: never break
inside a grapheme cluster, prefer to break at spaces, break
mid-word only if no break opportunity exists in the row.

\section{Nested clip}

When an element draws a child, the child's clip is the
intersection of its allocated rect and the parent's clip. So
a list inside a clipped panel never paints outside the panel
even if its rows logically extend beyond.

\begin{lstlisting}
Rect child_clip = clip_.intersect(child_rect);
\end{lstlisting}

The intersection is two min and two max operations.

\section{Scroll and clip}

A scrollable region has a viewport (the visible rectangle) and
a content size (often larger). The clip is the viewport; the
content is painted at a translated origin so the visible
portion lands inside the clip.

The trick: paint as if the content were at $(0, -\text{scroll}_y)$
relative to the viewport. Cells with $y < 0$ are clipped;
cells with $y \geq \text{viewport.h}$ are clipped. The middle
is shown.

\section{Clip and animation}

A sliding-in panel can be implemented with clipping. Draw the
panel at its final position; animate the parent's clip
rectangle to grow from empty to the full panel size. The
illusion is a wipe-in.

\maya{}'s reveal widget uses this pattern.

\section{Cell-perfect alignment}

A common bug: clip rectangle off by one. The right edge of a
border is included or excluded depending on convention. \maya{}
uses inclusive-exclusive: \code{Rect\{x, y, x+w, y+h\}} where
\code{x+w} is one past the last visible column.

This is the same convention as STL ranges. The arithmetic for
width is just \code{x2 - x1}.

\section{Debug overlay}

A debug build can enable a clip-rectangle overlay: every
element's clip is outlined in a thin tint. This is useful for
spotting layout bugs (an element drawing outside where you
thought it should).

\maya{} ships such an overlay behind a flag. Toggle with a
keybind.

\begin{tryit}
Enable the clip overlay in your app and toggle a panel that
shows/hides. Watch the rectangles update as layout reflows.
\end{tryit}

\section{Recap}

\begin{recap}
\begin{itemize}[leftmargin=*]
  \item Every element has a clip rectangle; writes outside
    are dropped.
  \item Five strategies for overflow: truncate, ellipsis,
    scroll, wrap, resize-parent.
  \item Nested clips intersect.
  \item Wide cells need careful truncation.
  \item Debug overlay helps diagnose layout bugs.
\end{itemize}
\end{recap}

\section*{Workshop}

\begin{enumerate}[leftmargin=*]
  \item Build a card with a long title. Compare truncate,
    ellipsis, and wrap. Which feels best?
  \item Resize the terminal small and confirm no element
    paints outside its rect.
  \item Implement a debug clip overlay; bind it to F12.
\end{enumerate}
